\documentclass[10pt,a4paper]{amsart}
\usepackage[margin=19mm]{geometry}
\usepackage{amsmath,amssymb,mathtools}
\usepackage[scr=rsfs,cal=euler]{mathalfa}
\usepackage{xfrac}
\usepackage[unicode,hidelinks,bookmarks=false]{hyperref}
\newcommand{\ie}{i.e.\ }
\newcommand{\cf}{cf.\ }
\newcommand{\resp}{respectively\ }
\newcommand{\Subsection}{\subsection}
\providecommand{\nobreakdash}{\nobreak\hbox{-}}
\setlength{\emergencystretch}{2em}
\multlinegap0pt
\usepackage{amssymb}
\usepackage{mathtools}
\usepackage{xcolor}
\usepackage[shortlabels]{enumitem}
\usepackage{tikz}
\usepackage[normalem]{ulem}
\newtheorem{cor}{Corollary}
\newtheorem{pf}{Pf}
\title{Convexity inequalities for eigenvalues and log-concavity of eigenfunctions}
\author{Paul Bryan, Julie Clutterbuck, Cale Rankin}
\date{Source: arXiv:2605.01334v1 (CC BY 4.0)}
\begin{document}
\maketitle

\noindent\emph{Source and licence.} Convexity inequalities for eigenvalues and log-concavity of eigenfunctions. Version arXiv:2605.01334v1 (CC BY 4.0), distributed by arXiv under the Creative Commons Attribution 4.0 International licence, http://creativecommons.org/licenses/by/4.0/, as recorded in the arXiv licence metadata for this exact version and shown on the abstract page https://arxiv.org/abs/2605.01334v1. Full licence notice: \texttt{licenses/arxiv-2605.01334-CC-BY-4.0.txt}.\par\smallskip

\noindent\textit{Editorial scope.} Counted unit: the article's cor org60a33f1 (source0.tex lines 320--329). The article's complete original proof of it is delivered with it (source0.tex lines 362--384, one \texttt{proof} environment of 23 lines, which the article places away from the statement). No mathematics is written, repaired or re-derived: the statement and the proof are the article's own lines, in the article's own notation, and no retained line is substituted or re-flowed.  Retained uncounted as the source-local context the proof needs: lines 335--340.  Nothing was omitted from any retained span, so the package carries no omission record.  No retained line contains a citation, so the wrapper carries no bibliography and no bibliography entry.  No retained line contains a figure environment, an image-inclusion command or drawing code, so no image is delivered and the package declares no asset dependency.  Editorial formatting only, in addition: an AMS \texttt{amsart} preamble that restates the article's own declarations and macros used by the retained lines, namely the environments \texttt{cor}, \texttt{pf} on separate counters, together with the standard packages those lines need.  The retained environments print in this document as Corollary 1, Pf 1 in order of appearance, whereas the article prints the same items with the numbering of its own full document.  The complete native source of the article as distributed in the arXiv e-print for this exact version is bundled unchanged as \texttt{sources/199630/source0.tex}.
\begin{cor}\label{eq:deriv-calc} Let $z \in \Omega_t$ and let $(x_0,x_1)$ be optimal
for $z$. If $\overline{u}$ is differentiable at $z$, then
\[ \frac{\nabla \overline{u}}{\overline{u}}(z) = (1-t)\frac{\nabla u_0}{u_0}(x_0) + t
\frac{\nabla u_1}{u_1}(x_1) = \frac{\nabla u_0}{u_0}(x_0) = \frac{\nabla u_1}{u_1}(x_1).
\]
\end{cor}

\begin{cor}
  \label{org60a33f1} Let\(u_i \in C^2(\Omega_i)\). Then the sup-convolution \(\overline{u}\) is locally Lipschitz
and Alexandrov twice differentiable almost everywhere. That is, there exists
measurable functions \(\nabla \overline{u} : \Omega_t \to \mathbf{R}^n\) and \(D^2 \overline{u} : \Omega_t
\to S_n\) (where \(S_n\) denotes \(n\times n\) nonnegative definite matrices) such that for almost
every \(z_0 \in \Omega_t\)
  \[ \overline{u}(z) = \overline{u}(z_0) + \langle \nabla \overline{u} (z_0), z-z_0\rangle +
\frac{1}{2} (z-z_0)^T D^2 \overline{u} (z_0) (z-z_0) + o (|z-z_0|^2)
  \] as \(z \to z_0\).
\end{cor}

\begin{pf} Since \(\Omega_t\) is bounded we have the following inclusions
  \[ \operatorname{Lip} \subseteq W^{1,\infty} \subseteq \dots W^{1,p} \subseteq \dots W^{1,1}.
  \]

From Lemma \ref{org60a33f1}, \(\overline{u}\) is locally Lipschitz. We just need
to provide a uniform Lipschitz constant to verify that \(\overline{u} \in
W^{1,p}\) for every \(p \geq 1\). Then since \(\overline{u}\) is continuous and
\(\overline{u} \equiv 0\) on the boundary \(\partial \Omega_t\) we will have \(\overline{u} \in
W^{1,p}_0\).

Since \(u_i\) are $C^1$ up to the boundary, \(|\nabla u_i|\) is uniformly bounded. Since \(\overline{u}\) is locally Lipschitz, it is differentiable almost
everywhere by Rademacher's theorem. At points where \(\overline{u}\) is
differentiable, Corollary \ref{eq:deriv-calc} ensures that
  \[ \frac{\nabla \overline{u}(z)}{u_0^{1-t}(x_0)u_1^t(x_1)} =\frac{\nabla
\overline{u}}{\overline{u}}(z) = \frac{\nabla u_0}{u_0}(x_0) = \frac{\nabla u_1}{u_1}
(x_1).
  \] In the case that \(u_1(x_1) \leq u_0(x_0)\) we have
  \[ |\nabla \overline{u}(z)| = \left|\frac{u_1^t(x_1)}{u_0^t(x_0)}\nabla u_0(x_0)\right|
\leq \left|\nabla u_0(x_0)\right| \leq \left\Vert \nabla u_0\right\Vert_{L^{\infty}}.
  \] Similarly, if \(u_0(x_0) \leq u_1(x_1)\) then $\left|\nabla \overline{u}(z)\right| \leq \left\|\nabla u_1\right\|_{L^{\infty}}$. Thus \(\overline{u}\) is locally Lipschitz with almost-everywhere uniformly
bounded gradient, hence is Lipschitz with Lipschitz constant bounded by \(\max
\{\|\nabla u_0\|_{L^{\infty}}, \|\nabla u_1\|_{L^{\infty}}\}\).
\end{pf}

\end{document}
